% !TEX root = ../main.tex

\chapter{Optimizations}
\label{chapter:optimizations}

% In order to perform optimizations on any code one must first identify the underlying bottlenecks and inefficiencies in that code.
% Hence, a simple profiler was added inside the code to measure the execution time of certain functions.
% The current chapter discusses how the optimizations were applied to the old code in order to achieve a speedup sufficient for the autonomous vehicle to respond in real time.
Before the DiffPlan stack can be optimized, the functions that dominate the runtime must first be identified.
The PyTorch \cite{paszke_PyTorchImperativeStyle_2019} profiler was added to the stack to record per-function execution times during simulation.
This chapter describes the two principal bottlenecks identified in the sampling planner and the loss calculation, explains how they were improved, and briefly discusses a subsequent GPU offloading attempt.
%! edit
The goal of each change is to bring the planner close to real-time execution on the onboard compute platform described in Chapter~\ref{chapter:Experiments}.

\section{Vectorize Sampling for Planner}
\label{section:VecSampling}

DiffPlan generates longitudinal candidate trajectories by solving a $3 \times 3$ boundary-value problem for each sampled terminal velocity $\dot{s}_{\text{end}, i}$.
The original implementation as abstracted in pseudocode is shown in Algorithm~\ref{alg:orig-long-sampling}, evaluating these trajectories inside nested Python loops.

\begin{algorithm}[ht!]
\caption{Original Longitudinal Sampling}
\label{alg:orig-long-sampling}
\begin{algorithmic}[1]
% \Require Initial position $s_{start}$, clamped initial velocity $\dot{s}_{init}$, time vector $\mathbf{t} \in \mathbb{R}^{m}$, target end velocities $\mathbf{v}_{end} \in \mathbb{R}^{N}$, duplication factor $n$, track length $L$
% \Ensure Position trajectories $\mathbf{S}$, velocity trajectories $\dot{\mathbf{S}}$, end positions $\mathbf{s}_{end}$
\State $t_{end} \gets \mathbf{t}[m]$ \Comment{Final time step}
\State $\mathbf{A} \gets \begin{pmatrix} 1 & 0 & 0 \\ 0 & 1 & 0 \\ 0 & 1 & 2 t_{end} \end{pmatrix}$
\State Initialize empty matrices $\mathbf{S}$, $\dot{\mathbf{S}}$ and empty vector $\mathbf{s}_{end}$
\For{$i = 1$ \textbf{to} $N$}
    \State $\mathbf{b} \gets \begin{pmatrix} s_{start} \\ \dot{s}_{init} \\ \mathbf{v}_{end, i} \end{pmatrix}$
    \State $\mathbf{c} \gets \mathbf{A}^{-1} \mathbf{b}$ \Comment{Solve for polynomial coefficients}
    \State $\mathbf{s}_{sample} \gets c_0 + c_1 \mathbf{t} + c_2 \mathbf{t}^2$ \Comment{Evaluate polynomial}
    \State $\dot{\mathbf{s}}_{sample} \gets c_1 + 2 c_2 \mathbf{t}$
    \State $\mathbf{s}_{sample} \gets \mathbf{s}_{sample} \ \text{mod} \ L$ \Comment{Track wrap-around}
    \State Append $\mathbf{s}_{sample, m}$ to $\mathbf{s}_{end}$ \Comment{Store final position}
    \For{$j = 1$ \textbf{to} $n$} \Comment{Duplicate samples}
        \State Append row vector $\mathbf{s}_{sample}^T$ to $\mathbf{S}$
        \State Append row vector $\dot{\mathbf{s}}_{sample}^T$ to $\dot{\mathbf{S}}$
    \EndFor
\EndFor
\State \Return $\mathbf{S}, \dot{\mathbf{S}}, \mathbf{s}_{end}$
\end{algorithmic}
\end{algorithm}

% \newpage
The outer loop runs $N$ times, once per sampled terminal velocity, and the inner loop duplicates each trajectory $n$ times before lateral sampling combines the results.
Although solving a single $3 \times 3$ system costs only $O(1)$ arithmetic, the implementation performs $N$ separate solves and $N \cdot n$ list append operations under the Python interpreter.
The overall complexity is therefore $O(N \cdot m \cdot n)$ in floating-point work, but the dominant cost is the $O(N \cdot n)$ sequential loop iterations.

The per-trajectory Python loop is removed by stacking all $N$ boundary vectors into one matrix $\mathbf{B} \in \mathbb{R}^{3 \times N}$ and solving $\mathbf{A}\mathbf{C}=\mathbf{B}$ in a single call.
Polynomial evaluation then becomes an outer product over the full time grid $\mathbf{t} \in \mathbb{R}^{m}$, producing all $N$ position and velocity profiles as dense tensors in $O(m \cdot N)$ time.
Algorithm~\ref{alg:vec-long-sampling} summarizes the vectorized procedure.

\begin{algorithm}[htbp]
\caption{Vectorized Longitudinal Sampling}
\label{alg:vec-long-sampling}
\begin{algorithmic}[1]
\Require Initial position $s_{start}$, clamped initial velocity $\dot{s}_{init}$, time vector $\mathbf{t} \in \mathbb{R}^{m}$, target end velocities $\mathbf{v}_{end} \in \mathbb{R}^{N}$, duplication factor $n$, track length $L$
\Ensure Position trajectories $\mathbf{S}$, velocity trajectories $\dot{\mathbf{S}}$, end positions $\mathbf{s}_{end}$
\State $t_{end} \gets \mathbf{t}[m]$ \Comment{Final time step}
\State $\mathbf{A} \gets \begin{pmatrix} 1 & 0 & 0 \\ 0 & 1 & 0 \\ 0 & 1 & 2 t_{end} \end{pmatrix}$
\State $\mathbf{B} \gets \begin{pmatrix} s_{start} & s_{start} & \dots & s_{start} \\ \dot{s}_{init} & \dot{s}_{init} & \dots & \dot{s}_{init} \\ \mathbf{v}_{end}[0] & \mathbf{v}_{end}[1] & \dots & \mathbf{v}_{end} [N-1] \end{pmatrix} \in \mathbb{R}^{3 \times N}$ 
\State $\mathbf{C} \gets \mathbf{A}^{-1} \mathbf{B}$ \Comment{Batched solve for all coefficients, $\mathbf{C} \in \mathbb{R}^{3 \times N}$}
\State Let $\mathbf{c}_0, \mathbf{c}_1, \mathbf{c}_2 \in \mathbb{R}^{1 \times N}$ be the row vectors of $\mathbf{C}$
\State Let $\mathbf{1}_{m} \in \mathbb{R}^{m}$ be a column vector of ones
\State $\mathbf{S}_{raw} \gets \mathbf{1}_m \mathbf{c}_0 + \mathbf{t} \mathbf{c}_1 + \mathbf{t}^2 \mathbf{c}_2$ \Comment{Outer product evaluation, $\mathbf{S}_{raw} \in \mathbb{R}^{m \times N}$}
\State $\dot{\mathbf{S}}_{raw} \gets \mathbf{1}_m \mathbf{c}_1 + 2\mathbf{t} \mathbf{c}_2$
\State $\mathbf{S}_{raw}^T \gets \mathbf{S}_{raw}^T \pmod{L}$ \Comment{Vectorized track wrap-around}
\State $\mathbf{s}_{end} \gets$ the final column vector of $\mathbf{S}_{raw}^T$
\State $\mathbf{S} \gets$ duplicate each row of $\mathbf{S}_{raw}^T$ sequentially $n$ times
\State $\dot{\mathbf{S}} \gets$ duplicate each row of $\dot{\mathbf{S}}_{raw}^T$ sequentially $n$ times
\State \Return $\mathbf{S}, \dot{\mathbf{S}}, \mathbf{s}_{end}$
\end{algorithmic}
\end{algorithm}

This reformulation reduces the number of sequential linear solves from $N$ to one, yielding an $N$-fold improvement in that phase.
In practice, this change will depend on the size of $N$ and the results for the runtime benchmarks can be found in Chapter~\ref{chapter:Results}.

\section{Vectorize Wall Proximity Calculation}
\label{section:VecWallCollision}

The profiler identified this function to take almost half the total execution time, which matches with the findings of Koch \cite{koch_DifferentiableInterpretablePath_2025}.
The collision loss depends on the distance between the ego vehicle and the nearest track-boundary point within a local search radius.
To approximate this distance, the planner discretizes the centerline into $N$ arc-length samples, where in this case $N = 50$ and computes left and right wall coordinates at each sample.
Algorithm~\ref{alg:orig-wall-collision} shows the original loop-based implementation.

\begin{algorithm}[htbp]
\caption{Original Wall Collision}
\label{alg:orig-wall-collision}
\begin{algorithmic}[1]
\Require Initial position $x_{car}$, $y_{car}$, search radius $r_{\text{search}} = 3.0$, number of discretized points $N = 50$
\Ensure Matrix of wall boundary points $\mathbf{P}_{wall} \in \mathbb{R}^{2N \times 2}$
\State Initialize empty list $\mathbf{P}_{\text{wall}}$
\State $s_{car} \gets \text{ProjectOnTrack}(x_{car}, y_{car})$ \Comment{Find car position in Frenet Frame}
\State $s_{\text{range}} \gets  1.5r_{\text{search}}$
\State $s_{\text{search}} \gets \text{LinSpace}(s_{car} - s_{\text{range}}, s_{car} + s_{\text{range}}, N)$ \Comment{$N$ evenly spaced points}
\For{$i = 1$ \textbf{to} $N$}
    \State $n_{\text{left}} \gets \text{TrackWidthLeft}(s_{\text{search}}[i])$
    \State $n_{\text{right}} \gets \text{TrackWidthRight}(s_{\text{search}}[i])$ \Comment{Evaluate track width}
    \State $p_{\text{center}} \gets \text{FrenetToCartesian}(s_{\text{search}}[i])$ \Comment{Centerline 2D coordinate}
    \State $\psi \gets \text{Heading}(s_{\text{search}}[i])$ \Comment{Heading of the car at s}
    \State $\mathbf{n} \gets \begin{pmatrix} -\sin\psi \\ \cos\psi \end{pmatrix}$ \Comment{Compute normal vector}
    \State $p_{\text{left}} \gets  p_{\text{center}} + n_{\text{left}} \cdot \mathbf{n}$
    \State $p_{\text{right}} \gets  p_{\text{center}} - n_{\text{right}} \cdot \mathbf{n}$
    \State Append $p_{\text{left}}$ and $ p_{\text{right}}$ to $\mathbf{P}_{\text{wall}}$
\EndFor
\end{algorithmic}
\end{algorithm}

Each loop iteration performs a track-width lookup, a Frenet-to-Cartesian conversion, a heading query, and two boundary-point updates.
The body is $O(1)$ per sample, but the loop runs $N$ times per collision evaluation, giving $O(N)$ total arithmetic with $N$ separate list appends.

The same logic is vectorized by evaluating every quantity across the full vector $s_{\text{search}}$ at once, as shown in Algorithm~\ref{alg:vec-wall-collision}.
Track-width queries, coordinate transforms, and heading lookups return tensors of shape $\mathbb{R}^{N}$ or $\mathbb{R}^{N \times 2}$.
Element-wise trigonometry forms the normal directions, and broadcasting computes all left and right boundary points in two matrix operations before concatenating the results into $\mathbf{P}_{\text{wall}}$.

\begin{algorithm}[htbp]
\caption{Vectorized Wall Collision}
\label{alg:vec-wall-collision}
\begin{algorithmic}[1]
\Require Initial position $x_{car}$, $y_{car}$, search radius $r_{\text{search}} = 3.0$, number of discretized points $N = 50$
\Ensure Matrix of wall boundary points $\mathbf{P}_{wall} \in \mathbb{R}^{2N \times 2}$
\State $s_{car} \gets \text{ProjectOnTrack}(x_{car}, y_{car})$ \Comment{Find car position in Frenet Frame}
\State $s_{\text{range}} \gets  1.5r_{\text{search}}$
\State $s_{\text{search}} \gets \text{LinSpace}(s_{car} - s_{\text{range}}, s_{car} + s_{\text{range}}, N)$ \Comment{$N$ evenly spaced points}
\State $\mathbf{n}_{\text{left}} \gets \text{TrackWidthLeft}(s_{\text{search}}) \in \mathbb{R}^N$ \Comment{Evaluate left widths for all points}
\State $\mathbf{n}_{\text{right}} \gets \text{TrackWidthRight}(s_{\text{search}}) \in \mathbb{R}^N$ \Comment{Evaluate right widths for all points}
\State $\mathbf{P}_{center} \gets \text{FrenetToCartesian}(s_{\text{search}}) \in \mathbb{R}^{N \times 2}$ \Comment{Centerline coordinates for all points}
\State $\boldsymbol{\psi} \gets \text{Heading}(s_{\text{search}}) \in \mathbb{R}^N$ \Comment{Headings for all points}
\State $\mathbf{n}_x \gets -\sin(\boldsymbol{\psi})$ \Comment{Element-wise sine operations}
\State $\mathbf{n}_y \gets \cos(\boldsymbol{\psi})$ \Comment{Element-wise cosine operations}
\State $\mathbf{n} \gets \begin{bmatrix} \mathbf{n}_x & \mathbf{n}_y \end{bmatrix} \in \mathbb{R}^{N \times 2}$ \Comment{Stack into normal matrix}
\State $\mathbf{P}_{\text{left}} \gets \mathbf{P}_{center} + \mathbf{n}_{\text{left}} \cdot \mathbf{n}$ \Comment{Scale normals (broadcast) and compute left points}
\State $\mathbf{P}_{\text{right}} \gets \mathbf{P}_{center} - \mathbf{n}_{\text{right}} \cdot \mathbf{n}$ \Comment{Scale normals (broadcast) and compute right points}
\State $\mathbf{P}_{wall} \gets \begin{bmatrix} \mathbf{P}_{\text{left}} \\ \mathbf{P}_{\text{right}} \end{bmatrix}$ \Comment{Concatenate points matrix vertically}
\State \Return $\mathbf{P}_{wall}$
\end{algorithmic}
\end{algorithm}

% \newpage
The arithmetic complexity remains $O(N)$ in the number of discretization points, but the number of Python loop iterations is reduced from $N$ to one.
This yields an $N$-fold reduction by allowing the use of SIMD parallelism \cite{flynn_ComputerOrganizationsTheir_1972, paszke_PyTorchImperativeStyle_2019}.
Because the collision loss is evaluated at every step, this vectorization removes the dominant bottleneck that previously limited the planner publish rate.

\section{GPU Offloading}
\label{section:GPUOffloading}

Considering that this code will eventually run on an NVIDIA Jetson Orin Nano Super, an embedded system-on-chip mounted on the RoboRacer vehicle, GPU offloading was explored as an additional optimization path.
The Jetson Orin Nano has a six-core Arm processor and an NVIDIA Ampere-equivalent GPU with 1024 CUDA cores and 32 tensor cores.
The Jetson provides substantial GPU compute and therefore should be explored by taking advantage of GPU offloading.
This is a common occurrence in neural networks, as they have an abundance of matrix-vector multiplications and other data-intensive operations for which the GPU is more suitable given its high core count compared to the CPU.
Hence, the central idea hinged on taking the vectorized code from Sections~\ref{section:VecSampling} and~\ref{section:VecWallCollision} and utilizing PyTorch's device context manager to transfer requested operations to the GPU.

Initially, the entire stack is offloaded to the GPU with the explicit goal of initializing all needed variables already on the GPU.
This required minimal effort to accomplish.
However, solving the LQR sub-problem of DiffMPC requires strict sequential ordering, since the solution at each time step depends on the previous one.
This caused the computation stack to be interrupted on the GPU, moving all variables and data to the CPU to compute the entire DiffMPC stack, for them to be moved back to GPU to continue the rest of the calculation.
The resulting host-device transfers overwhelmed the speedup expected from GPU offloading, producing adverse expected performance.

% when speaking about Felix and Jeremy's local planner we focus on the fact that the car rarely straifes away from the global planner but rather keeps going at that speed
% no advantages for the local portion of it other than obstacle avoidance but this is not very accurate to racing.